% !TEX root = ../main.tex
\chapter{Linear Rates: Products and Models}\label{ch:rates}
\begin{paracol}{2}

\begin{sources}
\begin{tabularx}{\linewidth}{@{}l l L@{}}
\toprule
\textbf{Lecture} & \textbf{Speaker} & \textbf{Content}\\
\midrule
2024 L7 & A.~Gunstensen (Mizuho) & Linear rates products and the size of their markets; LIBOR and its replacement by SOFR; discounting = funding; deposits, SOFR futures and interest-rate swaps; swap valuation and par swap rates; yield-curve construction (knot points, CDF and spline interpolation, calibration, FOMC curves, turns); delta ladders and gamma; P\&L attribution; clearing and the CCP basis; electronic trading and price making; Q\&A on the rates outlook, research roles and municipal bonds.\\
\bottomrule
\end{tabularx}

\smallskip
\textbf{Files.} Slides \emph{Linear Rates Models} \file{mit18_642_f24_lec07.pdf} (45 slides). Slide~39 (fast
model pricing) and slides 40--44 (bonds) were skipped in class for lack of time \ts{24}{7}{1:14:29}; they are
covered in \cref{ch06:sec-etrading,ch06:sec-bonds}. The data file \file{fedbondyielddata_20230911.csv} and the
notebook \file{feddataanalysis.ipynb} are listed on OCW together with Lecture~9 (PCA of Treasury yield changes)
and are treated in \cref{ch:pca}.
\textbf{Problem sets.} None; all exercises at the end are \emph{not from the course}.
\textbf{2013.} No counterpart: the 2013 course has no lecture on linear rates products. The dynamics of the
whole curve are the subject of the 2013 HJM lecture (\cref{ch:hjm}).
\end{sources}

\section*{Motivation}
Bonds, deposits, futures and swaps are the largest and most liquid markets in the world, and their payoffs are
fixed or linear in interest rates: pricing them requires no model of rate \emph{dynamics}, only today's
\emph{discount curve}. \Cref{ch:markets} took that curve as given (a flat rate, or yields read off a screen,
\cref{fig:yieldcurve}). This chapter asks the practitioner's questions: which rate should one discount with,
which instruments reveal the curve, how does one fill the gaps between a few dozen quotes, how does one hedge
and explain the daily P\&L of a book of swaps? The lecturer announced a deliberately ``non-mathy'' talk: a
house is designed by thinking about rooms and how they connect, not by spending most of the time on hammers
\ts{24}{7}{1:01}. The mathematics is elementary (linear algebra, interpolation, Taylor expansions); what matters
is the set of conventions and design choices, and the consequences of each choice for hedging. The curve built
here is the input of every rates-option model; its PCA is in \cref{ch:pca} and its stochastic evolution in
\cref{ch:hjm}.

% ------------------------------------------------------------------
\section{Linear products and their markets}\label{ch06:sec-markets}

\begin{definition}[Linear product]
A rates product is \emph{linear} if its payoff is either fixed or a linear function of observed interest rates
\ts{24}{7}{2:06}: (vanilla) bonds, certificates of deposit, interest-rate futures, forward rate agreements,
interest-rate swaps. There is no optionality (no caps, swaptions, callable or range-accrual features).
\end{definition}
The PV of a linear product is a linear function of discount factors and forward rates, so it is fixed by
today's curve alone: no volatility enters, except through small corrections (futures margining, averaging,
payment delays) discussed below. Why should a quant care \ts{24}{7}{3:07}?
\begin{itemize}
  \item \textbf{Size and liquidity} (slide 4). Bonds outstanding about \$128T (ICMA, 2020), interest-rate swaps
        \$573T notional (BIS, 2023), SOFR futures open interest above \$10T. Daily volumes: US Treasuries
        about \$900B, swaps \$800B, SOFR futures about 4 million contracts. In a market this size, a tiny margin
        is a large business \ts{24}{7}{3:07}.
  \item \textbf{Hedging and inputs.} Linear products hedge the complex ones, and their curves are inputs of the
        complex models: one builds up from the simple stuff \ts{24}{7}{4:08}.
  \item \textbf{Post-2008 complexity.} Before the crisis vanilla modelling was considered trivial. Funding,
        clearing and margining (\cref{ch06:sec-funding,ch06:sec-clearing}) made even ``\$1 paid in ten
        years'' hard to value, with no unique answer \ts{24}{7}{4:08}. Regulation (including model-risk management)
        ended much of the business in opaque, highly levered securitisations that had been sold to buyers who
        could not assess them \ts{24}{7}{5:08}.
\end{itemize}

% ------------------------------------------------------------------
\section{Interest rates and reference rates}\label{ch06:sec-rates}

\subsection{Simple interest and accrual}
An interest rate is the rate at which one earns (or pays) interest on a principal \ts{24}{7}{7:10}. Money-market
instruments use \emph{simple} interest:
\begin{equation}\label{ch06:eq-simple}
  I = N\,r\,\Delta t,\qquad
  \mathrm{PV}(X \text{ paid after } \Delta t) = \frac{X}{1+r_d\,\Delta t},
\end{equation}
where $\Delta t$ is the \emph{accrual fraction} of the period and $r_d$ a discounting rate. Example from class:
\$1000 in one year discounted at 5\% is worth $1000/1.05=\$952.38$ today \ts{24}{7}{8:11}. The interest rate
$r$ of an investment is easy to observe (the rate a bank advertises on a CD); the discounting rate $r_d$ is not,
and it depends on who is doing the discounting (\cref{ch06:sec-funding}).

\begin{coursenote}[Day counts (not discussed in class)]
$\Delta t$ is computed from dates with a \emph{day-count convention}: ACT/360 (actual days$/360$; SOFR, money
markets, the floating legs of USD swaps), ACT/365F, 30/360 (bond-style), ACT/ACT (Treasuries). Dates themselves
are moved off holidays by a \emph{business-day convention} (e.g.\ ``modified following''). The lecturer's advice
to job candidates: after a year on a swaps desk you will know \emph{date arithmetic}, and at least 90\% of all
problems with swaps are date problems \ts{24}{7}{34:51}. The number of days between two dates has no unique answer;
it depends on the convention.
\end{coursenote}

\subsection{LIBOR}
The \emph{London Interbank Offered Rate} was created in the late 1960s as a reference for syndicated bank
lending: corporate loans are mostly floating-rate, resetting every one or three months, and a common index was
needed \ts{24}{7}{9:13}. Every day a panel of about 20 banks reported the rate at which they could borrow
unsecured for various terms; the rate was a trimmed average of the submissions \ts{24}{7}{10:16}. The first
swaps (a cross-currency swap involving IBM in 1981, then single-currency interest-rate swaps) were built on it, and
by LIBOR's cessation some \$350--400T of swaps referenced it \ts{24}{7}{11:16}.

Its weaknesses surfaced after 2008 (slide 7):
\begin{itemize}
  \item it was a \emph{quoted}, not a transacted rate: submitters estimated where they \emph{would} borrow;
  \item it was \emph{unsecured}, and after Lehman almost nobody lent unsecured term money: some tenors saw a
        handful of real transactions in six months while a rate had to be published every day
        \ts{24}{7}{13:26};
  \item hence it could be manipulated. A book can have millions of dollars of P\&L per basis point
        (bp $=0.01\%$) of a fixing, and a group of traders coordinated with their banks' submitters to nudge the
        fixings; this led to prison sentences and multi-billion-dollar fines \ts{24}{7}{12:21}.
\end{itemize}
Regulators (led from London) required a replacement ``anchored in actual transactions and hard to
manipulate'' \ts{24}{7}{13:26}. USD LIBOR ended in June 2023; sterling earlier, Canada and Mexico around 2024
\ts{24}{7}{18:40}. The migration of mortgages, bonds, loans and hundreds of trillions of swaps was enormous
and tedious, comparable to the euro conversion and Y2K \ts{24}{7}{15:31}.

\begin{finance}[Why banks liked an unsecured index]
A bank's floating-rate loans pay LIBOR plus a fixed spread. When the banking system is under stress, bank
credit spreads widen and LIBOR rises, so the bank's \emph{income} rises exactly when its own funding becomes
expensive: a natural hedge (a ``right-way'' feature) \ts{24}{7}{17:40}, \ts{24}{7}{19:44}. A secured rate
contains no bank credit risk and lacks this property. The flip side: LIBOR-based hedges of instruments that
depend on government rates suffer when the LIBOR--Treasury spread blows out. In 2011--12 this cost the
lecturer's employer of the time a write-off of several billion dollars on structured notes (\$8.4B quoted in
class), most of which came back later \ts{24}{7}{14:31}.
\end{finance}

\subsection{SOFR}
A \emph{repo} (repurchase agreement) is a collateralised loan: I deliver \$100M of Treasuries and receive, say,
\$98M cash, agreeing to buy the bonds back the next day; the interest paid is the \emph{repo rate} (reverse
repo for the other side) \ts{24}{7}{16:35}. This is a real market with real volume (about \$800B a day
according to slide~8).

\begin{definition}[SOFR]
The \emph{Secured Overnight Financing Rate} is the main USD reference rate: a volume-weighted median of the
rates on overnight repo transactions collateralised by US Treasuries, published every business day
\ts{24}{7}{15:31}.
\end{definition}

\begin{coursenote}[Erratum (spoken)]
In class SOFR is expanded as ``Securitized Overnight Funding Rate'' and described as an average. The name is
\emph{Secured Overnight Financing Rate}, computed by the Federal Reserve Bank of New York (published since April
2018, selected in 2017 by the Fed-convened Alternative Reference Rates Committee) as a volume-weighted
\emph{median} of transaction rates. Similarly ``OAS rate'' in the transcript at 8:11 should read OIS
(overnight index swap) rate, as on slide~6.
\end{coursenote}

SOFR has three inconvenient features \ts{24}{7}{16:35}: it is a \emph{daily} rate, while nobody wants to pay
interest daily; a term rate built from it is known only at the \emph{end} of the period (set in arrears), while
borrowers like to know their payment in advance; and it is \emph{secured}, while banks prefer to lend
unsecured (previous box) \ts{24}{7}{17:40}.

\paragraph{From daily fixings to term rates} \ts{24}{7}{20:49}. Over an accrual period of $d$ days, let $r_i$
be the SOFR fixing of business day $i$, applying for $d_i$ calendar days ($d_i=3$ for a Friday fixing, which
covers the weekend), $d=\sum_{i=1}^nd_i$. The two standard term rates (slide 9) are
\begin{equation}\label{ch06:eq-sofr}
  R^{\mathrm{comp}} = \frac{360}{d}\Biggl[\,\prod_{i=1}^{n}\Bigl(1+\frac{r_i d_i}{360}\Bigr)-1\Biggr],
  \qquad
  R^{\mathrm{avg}} = \frac1d\sum_{i=1}^{n} r_i d_i .
\end{equation}
$R^{\mathrm{comp}}$ is exactly what one earns by rolling an overnight deposit (interest on interest), and
$R^{\mathrm{comp}}\ge R^{\mathrm{avg}}$ when rates are non-negative (\cref{ch06:ex-sofr}). Most USD products pay
compounded SOFR: it is simpler to calibrate, and it needs no correction for rate volatility, which an average
does \ts{24}{7}{21:49}. The reason is \cref{ch06:prop-float} below: a compounded leg is replicated exactly by
rolling a deposit.

\begin{example}[Compounding vs averaging (not from class)]
A 91-day period with a fixing every day ($d_i=1$): 5.30\% for 40 days, then 5.05\% after a 25\,bp FOMC cut.
Then $R^{\mathrm{avg}}=5.1599\%$ and $R^{\mathrm{comp}}=5.1933\%$: a 3.3\,bp difference. It is ``a few basis
points running'', irrelevant for a retail borrower but twenty times the bid--offer of a swap market maker.
\end{example}

\paragraph{Other rates} \ts{24}{7}{21:49}--\ts{24}{7}{26:16} (slide 10).
\begin{itemize}
  \item \textbf{Term SOFR}: 1-, 3-, 6- and 12-month rates published by the CME by fitting a model of daily rates to
        SOFR futures prices and averaging the fitted daily rates over the term \ts{24}{7}{24:02}. It is set in
        advance, like LIBOR, and popular in corporate loans, whose borrowers do not want to handle daily
        fixings. To prevent it from becoming a new LIBOR, the Fed restricts dealers to trading it only to hedge
        their loan books. The market is therefore one-way, and term-SOFR swaps trade 4--5\,bp away from vanilla
        SOFR swaps, 20--30 times the bid--offer \ts{24}{7}{25:09}.
  \item \textbf{Fed Funds}: the older unsecured overnight interbank rate, still widely used; plus Prime,
        DISCO (a Fannie Mae discount-note rate, the lecturer's favourite name), MMD (a municipal benchmark) and many
        others \ts{24}{7}{26:16}.
\end{itemize}

% ------------------------------------------------------------------
\section{Discounting equals funding}\label{ch06:sec-funding}

Which rate should discount a future cash flow \ts{24}{7}{26:16}? Assume (slides 11--12) we own a fixed payment
$F$ at time $T$, worth $P(t)$ today, and that we bought it with borrowed money: in this world nobody pays out
of pocket, positions are funded. Consider one day, from $t_0$ to $t_1$, with no market move
\ts{24}{7}{27:21}:
\begin{align*}
  \text{interest paid on the funding:}\quad & I(t_1) = P(t_0)\,\frac{r_b}{360},\\
  \text{accretion of the claim:}\quad & P(t_1) = P(t_0)\Bigl(1+\frac{r_d}{360}\Bigr),\quad\text{i.e. }\
  \Delta P = P(t_0)\,\frac{r_d}{360},
\end{align*}
where $r_b$ is our borrowing rate and $r_d$ the rate used to discount $F$.

\begin{proposition}[Discounting rate $=$ funding rate]\label{ch06:prop-funding}
Let an agent borrow and lend at the rate $r_b$ over $[t,T]$ ($n$ days, daily compounding, rates constant as on
the slides). The only price of $F$ paid at $T$ that gives the agent no riskless trade is
$P(t)=F\,(1+r_b/360)^{-n}$: future cash flows must be discounted at the funding rate, $r_d=r_b$.
\end{proposition}
\begin{proof}
The funded position has a daily P\&L $P(r_d-r_b)/360$, known in advance. If $r_d>r_b$, buying the claim with
borrowed money earns a riskless carry; if $r_d<r_b$, selling it and lending the proceeds does. Summing over the
holding period, the terminal P\&L of ``borrow $P$, buy the claim'' is $F-P(1+r_b/360)^n$, which vanishes for
every holding period only if $P$ is discounted at $r_b$ \ts{24}{7}{28:24}.
\end{proof}

\begin{remark}
If the agent borrows at $r_b$ but can only lend at $r_l<r_b$, the argument only gives the band
$F(1+r_b/360)^{-n}\le P\le F(1+r_l/360)^{-n}$. Different agents have different funding rates (the lecturer's
example: borrowers at 4\%, 6\% and 8\% depending on their credit) and therefore different values for the same
contract: there is no single price \ts{24}{7}{9:13}. The lecturer's one-sentence summary of the lecture:
\emph{funding and discounting are the same thing} \ts{24}{7}{28:24}.
\end{remark}

\begin{intuition}[A stationary state]
Think of the funded position as a closed system in which ``nothing happens'': the claim drifts up at $r_d$,
the loan at $r_b$. Zero P\&L in the absence of market moves (a stationary state) requires equal drifts. The
discount curve is not a law of nature but the rate at which the holder's money actually grows or costs;
compare \cref{sec:discounting}, where the same argument with a single universal rate gave $1/(1+r)$.
\end{intuition}

\begin{finance}[Funding desks and OIS discounting]
Banks do not let each trader manage the funding of a thousand counterparty relationships: a central funding
desk (treasury) lends to all desks at a generic internal rate and manages the funding bases and small
collateral exposures itself \ts{24}{7}{9:13}. For collateralised trades the relevant funding rate is the rate
paid on the collateral, an overnight rate: hence the market standard of discounting with an OIS (now SOFR)
curve \ts{24}{7}{8:11}. The costs of funding uncollateralised trades are priced as funding valuation
adjustments, part of the XVA family of \cref{ch:counterparty}.
\end{finance}

% ------------------------------------------------------------------
\section{The instruments that reveal the curve}\label{ch06:sec-instruments}
To build a curve we start from what trades actively and is easy to observe (slide 13): whatever the model, it
must reprice these instruments, because they can be traded \ts{24}{7}{28:24}. The workhorses are cash deposits
(CDs), interest-rate futures and swaps \ts{24}{7}{29:26}. We write each price as a function of the discount
function $Z(t)$ (value today of \$1 paid at $t$; $\Lambda_t$ in \cref{ch:markets}), then invert.

\subsection{Deposits and certificates of deposit}
Short-term money can be invested at a fixed simple rate $r(t)$ for maturities up to about a year (CDs,
interbank deposits; also repo and Treasury bills, not discussed in class). Investing 1 today returns
$1+r(t)\Delta$ at $t$, so (slide 14)
\begin{equation}\label{ch06:eq-deposit}
  1=\bigl[1+r(t)\Delta\bigr]Z(t)\qquad\Longrightarrow\qquad Z(t)=\frac{1}{1+r(t)\Delta}
  \quad\text{\ts{24}{7}{29:26}.}
\end{equation}
Deposits are well observed but awkward for a swaps desk: they are \emph{funded} (a \$100M CD requires \$100M of
cash and balance sheet), whereas derivatives need little cash up front and allow more leverage
\ts{24}{7}{30:32}.

\subsection{Forward rates and FRAs}
A floating payment that sets in the future can be replicated with fixed flows \ts{24}{7}{39:04}. Let the rate
apply to $[t_1,t_2]$ with accrual fraction $\Delta$, and suppose zero-coupon bonds maturing at $t_1$ and $t_2$
can be traded.

\begin{proposition}[Implied forward rate]\label{ch06:prop-fwd}
The fair (zero-PV) rate fixed today for lending $N$ from $t_1$ to $t_2$ is
\begin{equation}\label{ch06:eq-fwd}
  f(t_1,t_2)=\frac1\Delta\Bigl(\frac{Z(t_1)}{Z(t_2)}-1\Bigr).
\end{equation}
\end{proposition}
\begin{proof}
Lending $N$ at $t_1$ and receiving $N(1+f\Delta)$ at $t_2$ has present value
$-NZ(t_1)+Nf\Delta\,Z(t_2)+NZ(t_2)$. Setting it to zero gives $1+f\Delta=Z(t_1)/Z(t_2)$. If the market rate were
higher, lend forward and short the replicating zeros (buy the $t_1$ zero, sell the $t_2$ zero) for a riskless
profit, and vice versa.
\end{proof}
This is \cref{prop:forward} of \cref{ch:markets} with $F_{t_1t_2}=f\Delta$.

\begin{coursenote}[Erratum: slide 20]
The slide states the problem for a period $[t_1,t_2]$ but writes the equations with $t_0,t_1$:
``$\mathrm{PV}=-NZ(t_0)+Nf\Delta Z(t_1)+NZ(t_1)=0$'' and $f=\frac1\Delta\bigl(\frac{Z(t_0)}{Z(t_1)}-1\bigr)$; it also
says ``borrow $N$ at $t_0$''. The consistent statement is \eqref{ch06:eq-fwd}: the start of the accrual period
plays the role of $t_0$ on the slide, the end that of $t_1$.
\end{coursenote}

\paragraph{FRAs (not discussed in class).} A forward rate agreement on notional $N$ and strike $K$ pays at
the end of the period $N(R-K)\Delta$, where $R$ is the rate fixed for $[t_1,t_2]$ (in the LIBOR world
usually settled at $t_1$ as $N(R-K)\Delta/(1+R\Delta)$, the same value). By \cref{ch06:prop-fwd} its value today
is $N(f-K)\Delta\,Z(t_2)$ and the fair strike is $K=f$. A swap is a strip of FRAs sharing one strike.

\subsection{SOFR futures}\label{ch06:sec-futures}
Short-rate futures are among the most liquid contracts in every currency \ts{24}{7}{30:32}. The US contracts
trade at the CME (slide 15) \ts{24}{7}{31:35}:
\begin{itemize}
  \item \textbf{3-month SOFR futures} settle on the \emph{compounded} SOFR over a period running from one
        \emph{IMM date} to the next (third Wednesday of March, June, September, December; the dates come from
        the old International Monetary Market in Chicago);
  \item \textbf{1-month SOFR futures} settle on the \emph{average} SOFR over a calendar month.
\end{itemize}
Prices are quoted as $F=100-R$, with $R$ the settlement rate in percent: a 5\% rate gives 95, a 6\% rate 94
\ts{24}{7}{22:58}. The convention makes the contract behave like a bond: when rates fall (the market ``rallies''),
the price rises \ts{24}{7}{31:35}. Futures prices estimate forward rates up to a small \emph{convexity adjustment}
(slide 15):
\begin{equation}\label{ch06:eq-futures}
  f(t_1,t_2) = 100 - F - \mathrm{Cvx},\qquad \mathrm{Cvx}>0 .
\end{equation}

\paragraph{Origin of the adjustment} \ts{24}{7}{32:39}. Futures are margined daily: a contract entered at no
cost (apart from an initial margin) pays the holder the daily price change. A long position gains when the price
rises, i.e.\ when rates fall, and the gain can then be reinvested only at the lower rate; when rates rise, the
losses must be funded at the higher rate. This adverse correlation between the margin flows and the rate at
which they are carried makes a long futures position slightly worse than a long FRA, so the futures price is
lower and the futures rate $100-F$ higher than the forward rate. The adjustment can be estimated from a model, or
simply fitted together with swaps and futures prices.

\begin{remark}[Size of the adjustment (not discussed in class)]
Daily settlement makes the futures rate the expectation of $R$ under the risk-neutral (bank-account) measure,
$R^{\mathrm{fut}}=\E^{Q}[R]$, while the forward rate is its expectation under the $t_2$-forward measure, whose
density with respect to $Q$ is $D/\E^Q[D]$ with $D=\exp(-\int_0^{t_2}r_s\dd s)$ (change of numeraire,
\cref{ch:bs}). Hence
\[
  R^{\mathrm{fut}}-f=\E^Q[R]-\frac{\E^Q[DR]}{\E^Q[D]}=-\frac{\Cov^Q(D,R)}{\E^Q[D]}>0,
\]
because high rates come with low discount factors. In the Ho--Lee model with normal short-rate volatility
$\sigma$, $\mathrm{Cvx}\approx\tfrac12\sigma^2t_1t_2$ (continuous compounding; e.g.\ Hull's
\emph{Options, Futures, and Other Derivatives}). With $\sigma=1\%$ per year: 0.6\,bp for a contract one year
out, 13\,bp at five years, 51\,bp at ten years. It is negligible at the front and essential at the back, which
is why long-dated curves use swaps rather than futures.
\end{remark}

\subsection{Interest-rate swaps}
\begin{definition}[Vanilla interest-rate swap]
An agreement between two counterparties to exchange periodic cash flows on a notional $N$ until maturity
(slide 16) \ts{24}{7}{33:47}: the \emph{fixed leg} pays a rate $K$ set on the trade date, the \emph{floating leg}
pays a floating index observed periodically (today compounded SOFR). The \emph{payer} pays fixed and receives
floating; the \emph{receiver} does the opposite. Notionals are not exchanged.
\end{definition}
Swaps are the most liquid rates derivative. They convert fixed-rate into floating-rate exposure and vice
versa, are off balance sheet, cost little up front, and can be customised to any schedule a counterparty will
accept \ts{24}{7}{33:47}.

\begin{figure}[ht]
\centering
\begin{tikzpicture}[x=2.3cm,y=0.55cm,>=Stealth]
  \draw[->] (0,0) -- (5.55,0) node[right]{\small $t$ (years)};
  \foreach \k in {1,...,5} {
    \draw[->,thick,intcol] (\k,0) -- (\k,2.4);
    \node[below,font=\scriptsize] at (\k,-0.1) {$t_{\k}$};
  }
  \node[font=\scriptsize,intcol,anchor=south] at (5,2.4) {$NK\Delta_i$};
  \foreach \k/\h in {0.25/1.3,0.5/1.5,0.75/1.6,1/1.8,1.25/1.9,1.5/2.0,1.75/2.1,2/2.15,2.25/2.2,2.5/2.25,2.75/2.3,3/2.3,
                     3.25/2.35,3.5/2.4,3.75/2.45,4/2.5,4.25/2.55,4.5/2.6,4.75/2.65,5/2.7}
    \draw[->,thick,thmcol] (\k,0) -- (\k,-\h);
  \node[below left,font=\scriptsize] at (0,-0.1) {$t_0$};
  \node[font=\scriptsize,thmcol,anchor=north] at (5,-2.8) {$NR_j\Delta_j$};
  \node[font=\small,intcol,anchor=west] at (0.05,2.9) {fixed leg: received};
  \node[font=\small,thmcol,anchor=west] at (0.05,-3.3) {floating leg: paid, rate unknown today};
\end{tikzpicture}
\caption{Cash flows of a 5-year receiver swap with annual fixed payments $NK\Delta_i$ and a floating leg drawn
with quarterly payments $NR_j\Delta_j$ (USD SOFR swaps actually pay the compounded rate annually on both legs).
The floating amounts are unknown today; with an upward-sloping forward curve they are expected to grow.}
\label{ch06:fig-swapflows}
\end{figure}

\paragraph{Rolling out the cash flows.} The logic is simple but error-prone (slide 17) \ts{24}{7}{35:54}:
today $\to$ trade date; $+$ spot rule (typically two business days) $\to$ effective date; $+$ tenor $\to$
unadjusted maturity; step back by the payment frequency $\to$ unadjusted roll dates; $+$ business-day adjustment
$\to$ accrual dates; $+$ reset rule $\to$ rate-fixing dates; $+$ payment rule $\to$ payment dates. Every step
depends on holiday calendars, which change (Juneteenth, a state funeral, new national holidays), so models must
treat them as data. The standard calendar source lists 294 distinct dollar calendars (every exchange and business
centre has its own) \ts{24}{7}{37:00}. With the wrong calendar you misprice, and by the winner's curse you get
the trade only when the mispricing is against you.

% ------------------------------------------------------------------
\section{Swap valuation}\label{ch06:sec-swapval}
Let the fixed leg pay at $t_1<\dots<t_n$ with accrual fractions $\Delta_i$ and the floating leg at
$s_1<\dots<s_m$ with fractions $\Delta_j$, both starting at $t_0=s_0$.

\paragraph{Fixed leg} \ts{24}{7}{37:00}. Discount each known payment (slide 18):
\begin{equation}\label{ch06:eq-fixed}
  \mathrm{PV}_{\mathrm{fix}}=\sum_{i=1}^{n}K\Delta_iN\,Z(t_i)=NK\,A,\qquad
  A=\sum_{i=1}^{n}\Delta_iZ(t_i),
\end{equation}
where $A$ is the \emph{annuity} (or level, or PV01 of the fixed leg per unit rate).

\paragraph{Floating leg} \ts{24}{7}{38:04}. Replace unknown rates by forwards (slide 19),
$\mathrm{PV}_{\mathrm{float}}=\sum_{j}f_j\Delta_jN\,Z(s_j)$. This needs two curves: an \emph{index} (projection)
curve for the forwards $f_j$ and a \emph{discount} curve for the $Z(s_j)$. For vanilla SOFR swaps both are the
SOFR curve, so there is only one curve. For compounded SOFR the replication is exact:

\begin{proposition}[Floating leg]\label{ch06:prop-float}
Assume the discount curve is the SOFR curve (\cref{ch06:prop-funding} with SOFR as the funding rate) and that
each floating period pays $NR_j\Delta_j$ at $s_j$, where $R_j$ is SOFR compounded over $[s_{j-1},s_j]$ as in
\eqref{ch06:eq-sofr}. Then the value of period $j$ is $N\bigl[Z(s_{j-1})-Z(s_j)\bigr]=Nf_j\Delta_jZ(s_j)$ with
$f_j$ the forward rate \eqref{ch06:eq-fwd}, and
\begin{equation}\label{ch06:eq-float}
  \mathrm{PV}_{\mathrm{float}}=N\bigl[Z(s_0)-Z(s_m)\bigr].
\end{equation}
\end{proposition}
\begin{proof}
Invest $N$ at $s_{j-1}$ in an overnight deposit and roll principal and interest every day. At $s_j$ it has become
$N\prod_i(1+r_id_i/360)=N(1+R_j\Delta_j)$, whatever the fixings. Hence the payment equals ``value of the rolled
deposit at $s_j$'' minus $N$. The first part costs $N$ at $s_{j-1}$, worth $NZ(s_{j-1})$ today; the second is a
fixed amount worth $NZ(s_j)$. The equality with $Nf_j\Delta_jZ(s_j)$ is \eqref{ch06:eq-fwd}. Summing over $j$,
the sum telescopes.
\end{proof}

\begin{intuition}[A floating-rate note is worth par]
Adding the notional at the end, $\mathrm{PV}_{\mathrm{float}}+NZ(s_m)=NZ(s_0)$: a note paying the funding rate
is worth par at each reset, because its coupons are exactly what the money would have earned anyway. So a
receiver swap is ``long a fixed-coupon bond, short a floating-rate note'': a bond position financed at the
floating rate. This is also why an arithmetic \emph{average} of SOFR needs a correction: it is not what a rolled
deposit earns, so the static replication fails and a small, volatility-dependent adjustment appears
\ts{24}{7}{21:49}.
\end{intuition}

\begin{corollary}[Par swap rate]\label{ch06:cor-par}
A swap traded at zero PV (the market standard) has the \emph{par swap rate} \ts{24}{7}{40:06}
\begin{equation}\label{ch06:eq-par}
  S=\frac{\sum_{j=1}^mf_j\Delta_jZ(s_j)}{\sum_{i=1}^n\Delta_iZ(t_i)}
   =\frac{Z(t_0)-Z(t_n)}{A}.
\end{equation}
The first form (slide 21) is a weighted average of forward rates; the second holds under the single-curve
assumptions of \cref{ch06:prop-float}. For a spot-starting swap $Z(t_0)\approx1$.
\end{corollary}
Par swap rates are quoted on every broker screen, and they are what the curve must reproduce. With annual
payments, $\Delta_i=1$ and $Z(t_0)=1$, \eqref{ch06:eq-par} is exactly the par-bond equation used for
bootstrapping in \cref{sec:bonds}: the swap rate is the coupon of a par bond priced off the SOFR curve.

\paragraph{Mark-to-market and DV01.} After inception, let $S_t$ and $A_t$ be the par rate and annuity of the
remaining schedule. Then
\begin{equation}\label{ch06:eq-mtm}
  V_{\mathrm{payer}}=N\bigl[(Z(t_0)-Z(t_n))-KA\bigr]=N\,A_t\,(S_t-K),\qquad V_{\mathrm{receiver}}=-V_{\mathrm{payer}},
\end{equation}
and to first order (holding $A_t$ fixed) a 1\,bp rise in the par rate changes the payer's value by
$\mathrm{DV01}\approx N A_t\times10^{-4}$.

\begin{example}[Valuing swaps (not from class)]
With the curve of \cref{ch06:ex-boot} below: a 5-year receiver swap on \$100M struck at $K=4.00\%$, when the
5-year par rate is $3.40\%$ and $A_5=4.5178$, is worth $100\text{M}\times0.006\times4.5178=\$2.71$M. A new
10-year par swap on \$100M has $A_{10}=8.3195$ and a DV01 of about \$83{,}200 per bp.
\end{example}

\begin{finance}[Why swaps were invented, and why they trade at par]
A comment from the audience \ts{24}{7}{41:14}: in the late 1980s swaps boomed because highly rated companies
could borrow long-term fixed cheaply, while weaker companies could only borrow short-term floating. Each borrowed
where it had a comparative advantage and they swapped the payments, both ending up with cheaper funding of the
type they wanted \ts{24}{7}{42:14} (\cref{ch06:ex-compadv}). Today swaps are traded at (or within a bid--offer of
0.1--0.2\,bp around) zero PV, i.e.\ at the par rate. Off-market swaps appear inside structured transactions,
where the upfront value is accounted for elsewhere or is where the P\&L is taken \ts{24}{7}{41:14}.
\end{finance}

\begin{remark}[Two curves in the LIBOR era]
With LIBOR the index (unsecured, bank credit) and the discount rate (collateral, OIS) differed, so forwards came
from a LIBOR curve and discount factors from an OIS curve, and the telescoping \eqref{ch06:eq-float} failed
(``multi-curve'' pricing). The SOFR transition has largely brought vanilla USD swaps back to a single curve.
\end{remark}

% ------------------------------------------------------------------
\section{Building the yield curve}\label{ch06:sec-curve}
The yield curve is the core valuation model of a rates desk: a function $t\mapsto Z(t)$ that feeds all other
pricing \ts{24}{7}{42:14}. It is built in three steps (slide 22) \ts{24}{7}{43:15}:
\begin{enumerate}[1.]
  \item parameterise it by \emph{knot points} $\{(t_k,Z(t_k))\}_{k=1}^{N}$;
  \item choose an \emph{interpolation} to get $Z(t)$ for $t\notin\{t_k\}$ (say a 4.5-year discount factor from
        annual knots);
  \item \emph{calibrate} the $Z(t_k)$ to market quotes.
\end{enumerate}

\subsection{Knot points}
The simplest and most common choice: the maturities of the calibration instruments (1, 2, 3, 5, 7, 10-year swaps
give knots at 1, 2, 3, 5, 7, 10 years) \ts{24}{7}{43:15}. Other choices are possible (e.g.\ FOMC meeting dates,
\cref{ch06:sec-fomc}), but the knots must be chosen together with the interpolation and the instruments: bad choices
give unstable hedges \ts{24}{7}{44:16}. Asked why the knots are not chosen to minimise interpolation error (e.g.\
Chebyshev nodes), the lecturer answers that one wants a one-to-one link between knots and calibration instruments:
with knots at 5, 10, 15, \dots\ years and instruments concentrated at the short end, the curve would be
over-determined in one region and under-determined in another. Moreover the hedges come out in terms of the input
instruments, which must be liquid and tradeable \ts{24}{7}{45:19}.

\subsection{Interpolation: local versus global}
\begin{definition}[CDF interpolation]\label{ch06:def-cdf}
\emph{Constant daily forward} (CDF) interpolation assumes that the daily forward rate
$f_d$, defined by $Z(d+1)=Z(d)/(1+f_d/360)$, is constant between consecutive knots \ts{24}{7}{44:16}.
\end{definition}

\begin{proposition}[CDF $=$ log-linear discount factors]\label{ch06:prop-cdf}
Under CDF interpolation, for $t_k\le t\le t_{k+1}$,
\begin{equation}\label{ch06:eq-cdf}
  \ln Z(t)=\frac{t_{k+1}-t}{t_{k+1}-t_k}\ln Z(t_k)+\frac{t-t_k}{t_{k+1}-t_k}\ln Z(t_{k+1}),
\end{equation}
equivalently the product $t\,r_z(t)$ of time and continuously compounded zero rate $r_z(t)=-\ln Z(t)/t$ is
linear between knots (slide 23). The interpolation is \emph{local}: $Z(t)$ depends only on the two neighbouring
knots.
\end{proposition}
\begin{proof}
$\ln Z(d)=\ln Z(t_k)-\sum_{d'=t_k}^{d-1}\ln(1+f_{d'}/360)$ is linear in the number of days when $f_{d'}$ is
constant; matching the two endpoints gives \eqref{ch06:eq-cdf}, and $t\,r_z(t)=-\ln Z(t)$. In continuous time
this is a piecewise-constant instantaneous forward $f(t)=-\dd\ln Z/\dd t$.
\end{proof}
CDF is very fast, local, and produces the characteristic staircase forward curve (\cref{ch06:fig-fwd}): the
simplest, least smooth method \ts{24}{7}{46:21}.

\begin{coursenote}[Erratum: figure on slide 23]
The chart next to the CDF bullets is titled ``Yield'', but a step function is what CDF produces for the
\emph{forward} rate. The zero yield $r_z(t)=\frac1t\int_0^tf(s)\dd s$ is continuous (of the form $a+b/t$ on each
interval). Read the vertical axis as the daily forward rate.
\end{coursenote}

\begin{definition}[Cubic spline]
Fit on each interval $[t_k,t_{k+1}]$ a cubic
\[
  x(t)=\sum_{j=0}^{3}\alpha_{kj}\,(t-t_k)^j ,
\]
where $x$ is the interpolated quantity (typically the zero rate), and impose continuity of $x,x',x''$ at the
interior knots (slide 24) \ts{24}{7}{46:21}. Two end conditions (e.g.\ $x''=0$ at both ends, a \emph{natural}
spline) close the system: $4(N-1)$ coefficients for $2(N-1)$ interpolation conditions, $2(N-2)$ smoothness
conditions and 2 end conditions.
\end{definition}

\begin{coursenote}[Erratum: slide 24]
The slide writes $x_i(t)=\sum_{i=0}^{3}\alpha_{ij}(t-t_i)^i$, using $i$ both for the interval and for the
summation. The correct form is the one above: interval index $k$, power index $j$.
\end{coursenote}

The spline is reasonably fast and very smooth, but \emph{global}: the curve near an 8-year point depends on all
knots, because moving a distant knot changes the shape everywhere through the smoothness constraints
\ts{24}{7}{47:24}. The lecturer's advice to future rates quants is to keep a laminated card reading ``fast, local
hedges \emph{or} smooth, global hedges'': one can be anywhere in between, but not have both. He has argued about
this with traders for 30 years. Traders want an 8-year trade to be hedged with 7- and 10-year swaps only, not with
5- or 15-year ones. Among smooth methods the cubic spline is the worst in this respect, precisely because it is the
smoothest. Tension splines and B-splines sit in between, and much interpolation technology from computer graphics
could be used \ts{24}{7}{46:21} (slide 29).

\begin{intuition}[The spline as an elastic beam]
A natural cubic spline minimises the bending energy $\int x''(t)^2\dd t$ among all interpolants: it is the
draftsman's elastic ruler pinned at the knots. Moving one pin bends the whole ruler. For equally spaced knots the
disturbance decays by the factor $2-\sqrt3\approx0.27$ per knot \emph{with alternating sign}, the root of
$\lambda^2+4\lambda+1=0$ in the tridiagonal system for the second derivatives. This is exactly the ringing
visible in the spline hedges of \cref{ch06:fig-ladder}. CDF corresponds to a ``string'' of independent segments
in $\ln Z$: no stiffness, no propagation, no smoothness.
\end{intuition}

\subsection{Calibration}
For each calibration instrument write the pricing error as a function of the curve (slide 25)
\ts{24}{7}{48:27}:
\[
  E_i(Z)=P_i(Z)-P_i^{\mathrm{obs}},\qquad i=1,\dots,N .
\]
\begin{itemize}
  \item With $N$ knots and $N$ instruments solve $E_i=0$ exactly with any root finder. When the interpolation is
        local and instruments are sorted by maturity, each new instrument introduces a single new unknown and the
        system can be \emph{bootstrapped} one equation at a time (as in \cref{sec:bonds}).
  \item Otherwise minimise $E=\sum_iE_i(Z)^2$, possibly with extra terms, e.g.\ smoothness penalties on an
        under-specified curve, or accept any curve that prices every instrument within its bid--offer (which is
        indistinguishable from an exact fit) \ts{24}{7}{49:28}.
\end{itemize}

\begin{example}[Bootstrapping a small swap curve (not from class)]\label{ch06:ex-boot}
Annual-pay par swap rates ($\Delta_i=1$, single curve, day counts ignored): 1Y 4.00\%, 2Y 3.60\%, 3Y 3.45\%,
5Y 3.40\%, 7Y 3.45\%, 10Y 3.55\%. For $T\le3$ every payment date is a knot and \eqref{ch06:eq-par} gives the
recursion $Z_T=(1-S_T\sum_{k<T}Z_k)/(1+S_T)$. For 5Y, CDF gives $Z_4=\sqrt{Z_3Z_5}$ and $Z_5$ solves the scalar
equation $1-Z_5=S_5\bigl(Z_1+Z_2+Z_3+\sqrt{Z_3Z_5}+Z_5\bigr)$; similarly for 7Y and 10Y. Result:
\begin{center}\small
\begin{tabular}{@{}lcccccc@{}}
\toprule
knot $t_k$ (years) & 1 & 2 & 3 & 5 & 7 & 10\\
\midrule
par rate (\%) & 4.00 & 3.60 & 3.45 & 3.40 & 3.45 & 3.55\\
$Z(t_k)$ & 0.96154 & 0.93184 & 0.90351 & 0.84640 & 0.78874 & 0.70466\\
zero rate $r_z$ (\%, cont.) & 3.922 & 3.530 & 3.382 & 3.335 & 3.390 & 3.500\\
forward on $(t_{k-1},t_k]$ (\%) & 3.922 & 3.138 & 3.088 & 3.265 & 3.528 & 3.757\\
\bottomrule
\end{tabular}
\end{center}
A natural cubic spline on zero rates, calibrated to the same six quotes, reprices them exactly too, but it
implies a 4-year par rate of 3.405\% against 3.419\% for CDF. For a maturity that is not quoted, the choice of
interpolation moves the price by more than a basis point, about ten times the bid--offer.
\end{example}

\begin{figure}[ht]
\centering
\begin{tikzpicture}
\begin{axis}[width=0.8\linewidth,height=6.2cm,xlabel={maturity (years)},ylabel={rate (\%)},
  xmin=0,xmax=10.3,ymin=2.9,ymax=4.05,xtick={0,1,2,3,5,7,10},
  legend pos=south east,legend style={font=\small,cells={anchor=west}},grid=major,grid style={black!10}]
  \addplot[very thick,thmcol,const plot] coordinates {(0,3.922) (1,3.1375) (2,3.0875) (3,3.2649) (5,3.5277) (7,3.7574) (10,3.7574)};
  \addlegendentry{forward, CDF}
  \addplot[thick,defcol,smooth] coordinates {(1.00,3.471) (1.25,3.260) (1.50,3.093) (1.75,2.993) (2.00,2.981) (2.25,3.031)
    (2.50,3.093) (2.75,3.146) (3.00,3.174) (3.25,3.186) (3.50,3.204) (3.75,3.228) (4.00,3.257) (4.25,3.289) (4.50,3.323)
    (4.75,3.358) (5.00,3.394) (5.25,3.429) (5.50,3.465) (5.75,3.500) (6.00,3.534) (6.25,3.566) (6.50,3.595) (6.75,3.621)
    (7.00,3.643) (7.25,3.662) (7.50,3.682) (7.75,3.701) (8.00,3.721) (8.25,3.740) (8.50,3.760) (8.75,3.779) (9.00,3.798)
    (9.25,3.817) (9.50,3.836) (9.75,3.855) (10.00,3.874)};
  \addlegendentry{forward, cubic spline on zeros}
  \addplot[thick,black!55,dashed,mark=none] coordinates {(1,3.9221) (2,3.5298) (3,3.3824) (4,3.353) (5,3.3354) (6,3.3674)
    (7,3.3903) (8,3.4362) (9,3.4719) (10,3.5004)};
  \addlegendentry{zero rate, CDF}
  \addplot[only marks,mark=*,mark size=1.6pt,fincol] coordinates {(1,4.00) (2,3.60) (3,3.45) (5,3.40) (7,3.45) (10,3.55)};
  \addlegendentry{par swap quotes}
\end{axis}
\end{tikzpicture}
\caption{Curves implied by the quotes of \cref{ch06:ex-boot}. CDF gives a staircase of forward rates with jumps at
the knots; the spline (drawn from the first knot; below 1 year both methods coincide) gives a smooth forward curve
that dips below the CDF forwards around 2 years. Both reprice the six quotes exactly. Continuously compounded
rates; the 1-year zero rate is $\ln1.04=3.922\%$.}
\label{ch06:fig-fwd}
\end{figure}

\subsection{Central-bank meeting curves, turns and data-driven curves}\label{ch06:sec-fomc}
\paragraph{FOMC curves.} The history of daily SOFR (slide 26) is not smooth at all: it looks like a staircase
\ts{24}{7}{49:28}. It is flat between FOMC meetings and jumps when the Fed changes its policy rate: near zero
after March 2020, a sequence of steps up in 2022--23, a plateau around 5.3\% until the 50\,bp cut of September 2024
(``the chart needs a 50\,bp drop at the right edge'') \ts{24}{7}{50:32}. The spike at the left edge is a repo
squeeze related to tax-payment flows (September 2019). This suggests CDF interpolation with knots at the
meeting dates: daily rates are constant between meetings, which is now almost universal in the industry
\ts{24}{7}{50:32}. Two complications (slides 27--28):
\begin{itemize}
  \item there are 8 FOMC meetings a year but only 4 quarterly futures, so the problem is under-determined. One
        adds 1-month futures or smoothness constraints, or uses \emph{FOMC swaps} (meeting-to-meeting OIS), which
        now trade and directly price the rate expected after a given meeting \ts{24}{7}{51:34};
  \item the meeting calendar is known only about two years ahead, and instruments are sparse beyond the front end.
        Hence a \emph{hybrid} curve: CDF on meeting dates at the front, spline at the back, which works well in
        practice \ts{24}{7}{50:32}.
\end{itemize}

\paragraph{Turns} \ts{24}{7}{52:37}. Banks report balance sheets and return on assets at quarter- and year-ends and
want small balance sheets on those dates, so borrowing across year-end costs a premium: a spike in the forward
curve every year. Curves include these \emph{turns} explicitly. Some hedge funds trade pairs of swaps ending just
before and just after year-end, isolating the turn and exploiting whoever models it badly; in size this is real
money \ts{24}{7}{53:40}.

\paragraph{Data-driven curves} \ts{24}{7}{53:40}. A USD curve has about 50 inputs of very different liquidity;
perhaps half a dozen are super-liquid. One can estimate historical factors (curve shapes) by PCA
(\cref{ch:pca}), drive the first few factors from the most liquid instruments and build the curve from them. This
works well for bond curves (slide 29).

% ------------------------------------------------------------------
\section{Risk: PV01, delta ladders and gamma}\label{ch06:sec-risk}
On a desk, ``risk'' first means: how much do I make or lose if rates move \ts{24}{7}{54:42}?

\begin{definition}[PV01 and delta ladder]
The \emph{PV01} (or DV01) of a position is its sensitivity to a 1\,bp move of the curve inputs, computed by a
two-sided finite difference (slide 30)
\[
  \frac{\partial\mathrm{PV}}{\partial r}\approx\frac{\mathrm{PV}(r+\epsilon)-\mathrm{PV}(r-\epsilon)}{2\epsilon},
  \qquad \epsilon=1\,\text{bp}.
\]
Bumping each curve input $x_k$ separately gives the vector of partial PV01s, the \emph{delta ladder} (partial
deltas, bucketed risk, delta strip) $\delta_k=\partial\mathrm{PV}/\partial x_k$.
\end{definition}
The ladder is what the desk hedges and what risk limits are set on. It is computed by brute force (bump the input,
rebuild the curve, revalue) or analytically, but the concept is the same \ts{24}{7}{54:42}. If the inputs are the
calibration quotes, the ladder translates directly into hedges:

\begin{proposition}[Hedging with the calibration instruments]\label{ch06:prop-hedge}
Let the curve be calibrated exactly to par swaps with quotes $q=(q_1,\dots,q_N)$ and annuities $A_k$. A payer
swap on instrument $k$, struck at today's quote, has $\partial\mathrm{PV}_k/\partial q_j=A_k\,\delta_{kj}$ per unit
notional. Hence a portfolio with ladder $\delta_k=\partial V/\partial q_k$ is first-order neutral to every quote
after adding payer swaps with notionals
\[
  x_k=-\frac{\delta_k}{A_k},\qquad k=1,\dots,N .
\]
\end{proposition}
\begin{proof}
Because the calibrated curve reprices instrument $k$ at par rate $q_k$ for every $q$, \eqref{ch06:eq-mtm} gives
$\mathrm{PV}_k(q)=(q_k-K)A_k(q)$. Differentiating, $\partial_{q_j}\mathrm{PV}_k=\delta_{kj}A_k+(q_k-K)\partial_{q_j}A_k$,
and the second term vanishes at $K=q_k$. The total first-order sensitivity to $q_j$ is $\delta_j+x_jA_j=0$.
\end{proof}

\begin{example}[An 8-year swap hedged with 7- and 10-year swaps (not from class)]
Take the curve of \cref{ch06:ex-boot} and a \$100M 8-year receiver swap at its par rate (3.492\% with CDF). Its
parallel PV01 is $-\$68{,}835$ per bp, equal to $-NA_8\times10^{-4}$ with $A_8=6.883$. The ladder with respect to
the six quotes depends on the interpolation (\cref{ch06:fig-ladder}):
\begin{itemize}
  \item \textbf{CDF}: $-\$39{,}982$ on the 7Y quote, $-\$28{,}840$ on the 10Y, a few dollars elsewhere. The hedge is
        to pay fixed on \$65.3M 7Y and \$34.7M 10Y: exactly what traders want;
  \item \textbf{cubic spline}: $-533$, $+3{,}244$, $-6{,}969$, $+19{,}314$, $-68{,}766$, $-15{,}138$ dollars per bp
        on 1Y, \dots, 10Y. The hedge requires positions of alternating sign in every maturity, including
        \$112M of 7Y and $-\$43$M of 5Y: a ``smeared hedge'' (slide 32).
\end{itemize}
In both cases the ladder sums to the parallel PV01.
\end{example}

\begin{figure}[ht]
\centering
\begin{tikzpicture}
\begin{axis}[width=0.85\linewidth,height=6cm,ybar,bar width=9pt,
  symbolic x coords={1Y,2Y,3Y,5Y,7Y,10Y},xtick=data,
  ylabel={hedge: payer notional (\$M)},xlabel={calibration swap},
  ymin=-55,ymax=125,legend pos=north west,legend style={font=\small},
  ymajorgrids,grid style={black!10},
  nodes near coords,every node near coord/.append style={font=\tiny},
  point meta=explicit symbolic]
  \addplot[fill=thmcol!70,draw=thmcol] coordinates {(1Y,0) [0] (2Y,0) [0] (3Y,0) [0] (5Y,0) [0] (7Y,65.3) [65.3] (10Y,34.7) [34.7]};
  \addlegendentry{CDF (local)}
  \addplot[fill=defcol!60,draw=defcol] coordinates {(1Y,5.5) [5.5] (2Y,-17.1) [$-$17.1] (3Y,24.9) [24.9] (5Y,-42.8) [$-$42.8] (7Y,112.3) [112.3] (10Y,18.2) [18.2]};
  \addlegendentry{cubic spline (global)}
\end{axis}
\end{tikzpicture}
\caption{Hedge of a \$100M 8-year receiver swap with the six calibration swaps, $x_k=-\delta_k/A_k$ from
\cref{ch06:prop-hedge}, for the two interpolations of \cref{ch06:ex-boot}. Both hedges total about \$100M, but
the spline one rings with alternating signs through the whole curve (``fast, local'' versus ``smooth, global''
\ts{24}{7}{47:24}).}
\label{ch06:fig-ladder}
\end{figure}

\paragraph{Gamma} \ts{24}{7}{55:47}. Second-order risk (gamma, or convexity) matters more or less depending on the
book. Slide 31 defines it as the change in delta for a 1\,bp move:
\[
  \Gamma=\frac{\mathrm{PV}(r+\epsilon)-2\,\mathrm{PV}(r)+\mathrm{PV}(r-\epsilon)}{\epsilon}
  =\epsilon\,\frac{\partial^2\mathrm{PV}}{\partial r^2}+O(\epsilon^3).
\]
Note the single $\epsilon$ in the denominator: this is $\Delta(r+\tfrac\epsilon2)-\Delta(r-\tfrac\epsilon2)$ with
$\Delta=\partial\mathrm{PV}/\partial r$, i.e.\ the second derivative times the bump, not the second derivative
itself. The full matrix of cross gammas $\partial^2\mathrm{PV}/\partial x_j\partial x_k$ needs bumps in pairs, of
order $N^2$ revaluations: $2{,}500$ for a 50-input curve, too expensive. Desks use a parallel gamma, a factor (PCA)
gamma, or scenario analysis instead \ts{24}{7}{55:47}.

\paragraph{Practical issues} \ts{24}{7}{56:49} (slide 32).
\begin{itemize}
  \item \emph{Bump size}: 1\,bp is a reporting convention. Monte Carlo pricers need larger bumps, or the noise
        swamps the difference (the ``noisy delta'' problem of \cref{ch:markets}); PDE pricers show grid effects
        for large bumps; too large a bump creates unnatural curve shapes.
  \item \emph{Sum of partials vs parallel shift}: by the chain rule they agree to first order, as in the example;
        the discrepancy comes from cross-gamma terms at finite $\epsilon$. Sequential bumps are one remedy.
  \item \emph{Interpolation}: the ladder of a position is a property of position \emph{and} curve model. Local
        interpolation concentrates it near the maturity, smooth interpolation smears it, which is a source of
        endless arguments \ts{24}{7}{56:49}.
\end{itemize}

% ------------------------------------------------------------------
\section{P\&L attribution}\label{ch06:sec-pnl}
Every morning a desk starts with positions and a market, and every evening it has made or lost money: why
\ts{24}{7}{56:49}? Skill, market direction, client flow earning the bid--offer? The attribution is also one of the
best checks of the model, and in some jurisdictions a regulatory requirement. Total P\&L is known exactly (slide 33),
$\mathrm{PnL}(t_0,t_1)=\mathrm{PV}(t_1)-\mathrm{PV}(t_0)$. It is split into three parts (slide 34)
\ts{24}{7}{57:58}:
\begin{itemize}
  \item \textbf{Portfolio change}: new trades, unwinds, novations, resets, exercises. Conceptually simple,
        mechanically the hardest (booking systems, operational flows).
  \item \textbf{Time}, or \emph{theta}: the change in value if the market ``does not move''.
  \item \textbf{Market moves}: to first order $\sum_i\frac{\partial\mathrm{PV}}{\partial x_i}\Delta x_i$ (the
        linear term on slide 34), plus $\frac12\sum_{i,j}\Gamma_{ij}\Delta x_i\Delta x_j$ and higher terms.
\end{itemize}
What remains is the \emph{unexplained} P\&L. A large or systematic unexplained part signals missing risk factors or
a bad model. The order in which the three steps are applied matters (the decomposition is path-dependent), so
conventions must be fixed and documented.

\paragraph{What does ``the market did not move'' mean?} \ts{24}{7}{58:59}. The curve is a function of maturity,
and when the calendar moves one must decide what is held fixed. There is more than one defensible answer
(not detailed in class):
\begin{itemize}
  \item \emph{Forwards realised}: yesterday's forward curve for fixed calendar dates is today's curve,
        $Z_{t_1}(T)=Z_{t_0}(T)/Z_{t_0}(t_1)$. Every funded position then earns exactly the funding rate: theta
        is pure carry.
  \item \emph{Constant curve by tenor}: the zero (or par) rate for each time-to-maturity is unchanged. On an
        upward-sloping curve a bond ``rolls down'' to lower yields and earns carry \emph{plus} roll-down
        (\cref{ch06:ex-theta}).
\end{itemize}
For a physicist: the first convention is Lagrangian (follow each cash flow), the second Eulerian (sit at a fixed
point of the maturity axis). Whatever the choice, the desk and its controllers must use the same one.

\begin{coursenote}[The spoken P\&L example]
In class the market term is illustrated with a long position of \$1B of bonds, a PV01 of ``\$10 million'' and
gains of \$50M for a 5\,bp move \ts{24}{7}{58:59}. The numbers were improvised: \$1B of 10-year bonds has a PV01
of roughly \$0.8--0.9M, so 5\,bp is worth about \$4M, and a long bond position gains when rates \emph{fall}.
\end{coursenote}

% ------------------------------------------------------------------
\section{Clearing, margin and the CCP basis}\label{ch06:sec-clearing}
A swap starts near zero PV, but as markets move its value drifts away from zero and becomes a credit exposure: if
a counterparty that owes me \$10M on a swap defaults, I lose it (slide 35) \ts{24}{7}{1:00:00}. Lehman, a huge
swap dealer, made this concrete in 2008: its counterparties had to replace hundreds of billions of hedges at once,
bid--offers widened about twentyfold for a month or two, and the one large dealer still trading normally captured
the flow \ts{24}{7}{1:01:03}. After the crisis regulators made \emph{central clearing} mandatory for vanilla
swaps: once two parties agree a trade, both assign their side to a central counterparty (CCP: LCH in London,
CME, JSCC in Japan) and post margin to it, so they no longer face each other's credit \ts{24}{7}{1:01:03}.

\begin{finance}[Fewer, larger risks]
Clearing replaces many small bilateral risks by concentrated exposure to a few enormous, deeply interconnected
CCPs. A failure of one of them would be a systemic catastrophe, which is why members post margin and default-fund
contributions and CCPs are heavily capitalised \ts{24}{7}{1:02:09}.
\end{finance}

\paragraph{The CCP basis} \ts{24}{7}{1:02:09}. A dollar owed by LCH is not the same as a dollar owed by CME:
positions at different CCPs cannot be offset operationally or for collateral. From about 2017 the fixed rate for
paying fixed at CME exceeded that at LCH. On a broker screen for 26 June 2017 (slide 36; source Tradition) the
CME--LCH spread grew with maturity, from 0.10\,bp at 1 year through 1.05 (5Y), 2.45 (10Y) and 3.40 (30Y) to
3.55\,bp at 50 years, in a market with a bid--offer of about 0.2\,bp \ts{24}{7}{1:03:10}. It persisted for years
(recently it has changed sign). Why not receive at CME, pay at LCH and retire \ts{24}{7}{1:03:10}?
\begin{itemize}
  \item \textbf{Double margin}: the two offsetting trades sit at different CCPs, so initial margin (1--2\% of
        notional, quoted in class) is posted twice. Its funding cost is the main driver
        \ts{24}{7}{1:04:16}.
  \item \textbf{Netting}: CME also clears futures, so a dealer long futures and short swaps there gets partial
        margin offsets that the arbitrage would destroy \ts{24}{7}{1:04:16}.
  \item \textbf{Preferences}: large US asset managers, one-way payers of fixed, preferred CME, a US-jurisdiction
        CCP not controlled by dealers, over LCH \ts{24}{7}{1:05:21}. (LCH Group is today majority-owned by the
        London Stock Exchange Group, with banks as minority shareholders; the perception quoted in class concerns
        its governance.)
\end{itemize}
The general lesson: margin and capital requirements let spreads persist between instruments that would otherwise
be arbitraged. To price correctly one must include how the trade is margined, how much capital it consumes, and
every other cost that used to be ignored. Estimating these costs and pushing them into the upfront price has been
one of the largest quant efforts of the last 15 years, made harder by the fact that many are portfolio effects
that must be allocated to individual trades \ts{24}{7}{1:06:21} (see \cref{ch:counterparty}).

% ------------------------------------------------------------------
\section{Electronic trading and price making}\label{ch06:sec-etrading}
Electronic trading is one of the hottest areas for quants \ts{24}{7}{1:07:23}. As a rule of thumb, equities are
about ten years ahead of FX, and FX ten years ahead of fixed income. Slide 37 (Flow Traders data, 2023) shows the
share of electronic trading: above 90\% for US cash equities and equity options, European ETFs and FX spot;
about 80\% for European government bonds and 65\% for US Treasuries; 35--45\% for investment-grade corporates;
22--27\% for high yield; about 12\% for munis and 10\% for emerging-market bonds. Small trades go electronic first:
\$1M of Treasuries does not move a market trading \$1T a day, whereas \$500M does and needs care
\ts{24}{7}{1:08:24}.

\paragraph{Participants} \ts{24}{7}{1:10:26}. \emph{Principal trading firms} (high-frequency shops) have no
customers, trade their own capital and are extremely fast; the lecturer met a physics PhD hired to optimise the
microwave links between Chicago and New York, and realised that banks will never beat them at latency
\ts{24}{7}{1:11:26}. \emph{Dealers} (banks) are fast (microseconds, not nanoseconds) and have customers, who pay
the bid--offer. \emph{Customers} want efficient execution, a view of the full market and confidence that they are
not being taken advantage of \ts{24}{7}{1:12:26}.

\paragraph{Price making} (slide 38) \ts{24}{7}{1:12:26}. Start from a mid price, from liquidity providers for
actively traded instruments or from a model for off-the-run ones, then \emph{skew} it according to:
\begin{itemize}
  \item the current position: be aggressive on risk-reducing trades, conservative on risk-increasing ones;
  \item the client tier: a valuable client relationship justifies a better price \ts{24}{7}{1:13:29};
  \item \emph{client toxicity}: the classic toxic client sends the same request to ten dealers at once. Whoever
        wins finds nine others hedging the same way, and the hedge has moved away (winner's curse)
        \ts{24}{7}{1:13:29};
  \item market volatility (widen when the market may move before the hedge is done) and short-term prediction
        models (market microstructure).
\end{itemize}
All of this requires tightly integrated systems. A credit check that answers after the client has gone is a lost
trade \ts{24}{7}{1:14:29}.

\paragraph{Fast model pricing (slide 39, not presented in class).} Model prices must be recomputed in microseconds.
One linearises the output around the last full valuation,
\[
  r(t)\approx r(t_0)+\sum_{i=1}^{n}\frac{\partial r}{\partial x_i}\,\Delta x_i ,
\]
splits the drivers into fast ones (fed through the linearisation in real time) and slow ones, and periodically
recomputes the full model. The slide omits the index on $\Delta x$; it must be $\Delta x_i$.

% ------------------------------------------------------------------
\section{Bonds revisited}\label{ch06:sec-bonds}
Slides 40--44 were not presented in class; they extend \cref{sec:bonds,sec:duration}.

\paragraph{Why bonds are harder than swaps.} Government bonds are the most liquid fixed-income instruments in many
markets. G20 government bonds are usually treated as default-free, while emerging-market bonds need a default model.
Despite their apparent simplicity bonds are harder than swaps: bonds are not fungible (each issue has its own
coupon, maturity, liquidity, repo specialness), whereas all swaps are generated from one curve (slide 40).

\paragraph{Price--yield with $f$ payments a year} (slide 41). Bonds trade and settle on price, but yield is the
better comparison across coupons and maturities. With $n$ remaining periods, coupon $C_i$ per period per unit
notional and $f$ coupons per year,
\[
  P_D(y)=\sum_{i=1}^{n}\frac{C_i}{(1+y/f)^{i}}+\frac{1}{(1+y/f)^{n}},
\]
which reduces to \eqref{eq:couponbond} for $f=1$. The slide's chart shows the familiar picture: premium when the
yield is below the coupon, discount when above.

\begin{proposition}[Modified and Macaulay duration]
Let $D_{\mathrm{mod}}=-\frac1P\frac{\partial P}{\partial y}$ and let $D_{\mathrm{Mac}}=\sum_i\frac{i}{f}\,
\frac{\mathrm{CF}_i(1+y/f)^{-i}}{P}$ be the PV-weighted average payment time in years. Then
$D_{\mathrm{Mac}}=D_{\mathrm{mod}}\,(1+y/f)$ (slide 42).
\end{proposition}
\begin{proof}
$\partial_y(1+y/f)^{-i}=-\frac{i}{f}(1+y/f)^{-i-1}$, so
$-P'(y)=(1+y/f)^{-1}\sum_i\frac if\mathrm{CF}_i(1+y/f)^{-i}=P\,D_{\mathrm{Mac}}/(1+y/f)$.
\end{proof}
(The slide writes the factor as $1+y/n$ with $n$ the compounding frequency, called $f$ on the previous slide.)
Duration is ``not generic but good to gain insight'': it assumes a flat move of a single yield, while swap risk is
a ladder.

\paragraph{Bonds against the swap curve: the z-spread} (slide 43). Swap and bond books hedge each other, so their
risks must be aggregated. Pricing a bond off the swap (SOFR) curve,
$P'_D=\sum_iC_iZ(t_i)+Z(t_n)$, will not reproduce its market price; the \emph{z-spread} $s$ is the parallel shift
of the continuously compounded discount rates that does:
\[
  P_D=\sum_{i=1}^{n}C_iZ(t_i)e^{-st_i}+Z(t_n)e^{-st_n}.
\]
It measures the bond's credit, liquidity and supply premium relative to swaps, and it lets the bond's rate risk be
expressed as a delta ladder on the swap curve plus a single spread risk.

\begin{example}[z-spread and DV01 (not from class)]
A 5-year bond with 4\% annual coupons trades at 101.50. Off the curve of \cref{ch06:ex-boot} it would be worth
$P'_D=102.71$; the z-spread is $s=25.6$\,bp. Its yield is 3.666\%, $D_{\mathrm{Mac}}=4.633$,
$D_{\mathrm{mod}}=4.469$, and its DV01 is $D_{\mathrm{mod}}\,P\times10^{-4}=0.0454$ per 100 of face value.
\end{example}

\paragraph{Bond flavours} (slide 44): government, asset-backed, corporate, municipal, high yield and more, each with
its own valuation and risk nuances. \emph{Municipal bonds} were discussed in class \ts{24}{7}{1:08:24}: a huge number
of small and heterogeneous issuers (states, towns, school boards, toll-road and bridge projects) and interest exempt
from federal income tax. At a tax rate $\tau$ an investor is indifferent between a muni yielding $y_m$ and a taxable
bond yielding $y_T$ when $y_m=(1-\tau)y_T$, so munis yield about 70\% of comparable taxable bonds
\ts{24}{7}{1:09:26}. Only taxpayers benefit, so the market is retail: the lecturer's former electronic muni platform
(acquired by TD) did 5{,}000--10{,}000 trades a day with an average size around \$28{,}000, absurdly small by
government-bond standards but profitable in volume. Emerging-market bonds are idiosyncratic (payment terms, credit
support) and still mostly traded by phone \ts{24}{7}{1:10:26}.

\begin{finance}[From the Q\&A]
\emph{Rates outlook} \ts{24}{7}{1:15:32}. The Fed balances inflation, pushed up by COVID-era stimulus and by the
energy shock after Russia's invasion of Ukraine, against employment. High rates discourage borrowing (business
start-ups, and above all housing, the largest purchase in most lives), and a cutting cycle had just begun.
\emph{Two kinds of fixed-income research} \ts{24}{7}{1:16:39}: client-facing economics research makes forecasts
and trade recommendations; internal quant research on pricing and hedging makes no predictions at all. It only
needs to know how to hedge (``if I knew where rates are going, I would be on my Caribbean island'').
\emph{Muni analytics} \ts{24}{7}{1:17:42}: revenue bonds depend on project cash flows (toll roads lost their revenue
in COVID); all munis depend on future tax policy, since the value of the exemption moves with the tax rate
(\cref{ch06:ex-muni}); and the practical problem is \emph{sourcing}: finding who holds the bonds a client wants,
which is why munis are still voice-driven \ts{24}{7}{1:18:43}.
\end{finance}

% ------------------------------------------------------------------
\begin{keyformulas}
\begin{align*}
&\text{Simple interest, deposit:} && I=Nr\Delta t,\\
& && Z(t)=1/\bigl(1+r(t)\Delta\bigr)\\
&\text{SOFR term rates:} && R^{\mathrm{comp}}=\tfrac{360}{d}\bigl[\textstyle\prod_i(1+\tfrac{r_id_i}{360})-1\bigr],\\
& && R^{\mathrm{avg}}=\tfrac1d\sum_ir_id_i\\
&\text{Discounting:} && r_d=r_b\ \text{(discount at the funding rate)}\\
&\text{Forward rate:} && f(t_1,t_2)=\tfrac1\Delta\bigl(Z(t_1)/Z(t_2)-1\bigr)\\
&\text{Futures:} && f=100-F-\mathrm{Cvx},\\
& && \mathrm{Cvx}\approx\tfrac12\sigma^2t_1t_2\ \text{(Ho--Lee)}\\
&\text{Swap legs:} && \mathrm{PV}_{\mathrm{fix}}=NK\textstyle\sum_i\Delta_iZ(t_i)=NKA,\\
& && \mathrm{PV}_{\mathrm{float}}=N\bigl[Z(t_0)-Z(t_n)\bigr]\\
&\text{Par rate, MTM:} && S=\bigl(Z(t_0)-Z(t_n)\bigr)/A,\\
& && V_{\mathrm{payer}}=NA(S-K),\quad \mathrm{DV01}\approx NA\cdot10^{-4}\\
&\text{CDF interpolation:} && \ln Z\ \text{linear between knots}\iff\\
& && f(t)\ \text{piecewise constant}\\
&\text{Calibration:} && E_i=P_i(Z)-P_i^{\mathrm{obs}}=0\quad\text{or}\quad\min\textstyle\sum_iE_i^2\\
&\text{PV01, gamma:} && \tfrac{\mathrm{PV}(r+\epsilon)-\mathrm{PV}(r-\epsilon)}{2\epsilon},\\
& && \Gamma=\tfrac{\mathrm{PV}(r+\epsilon)-2\mathrm{PV}(r)+\mathrm{PV}(r-\epsilon)}{\epsilon}\\
&\text{Hedge with inputs:} && x_k=-\delta_k/A_k\\
&\text{P\&L explain:} && \mathrm{PnL}=\text{portfolio}+\Theta\,\Delta t\\
& && +\textstyle\sum_i\delta_i\Delta x_i+\tfrac12\sum_{ij}\Gamma_{ij}\Delta x_i\Delta x_j+\text{unexplained}\\
&\text{Bonds:} && D_{\mathrm{Mac}}=D_{\mathrm{mod}}(1+y/f),\\
& && P_D=\textstyle\sum_iC_iZ(t_i)e^{-st_i}+Z(t_n)e^{-st_n}\\
&\text{Munis:} && y_m=(1-\tau)\,y_T
\end{align*}
\end{keyformulas}

% ------------------------------------------------------------------
\section*{Exercises}
\addcontentsline{toc}{section}{Exercises}
No problem set covers this lecture; all exercises are \emph{not from the course}.

\begin{exercise}[not from the course]\label{ch06:ex-sofr}
\leavevmode
\begin{enumerate}[(a)]
  \item Show that $R^{\mathrm{comp}}\ge R^{\mathrm{avg}}$ in \eqref{ch06:eq-sofr} whenever all $r_i\ge0$.
  \item For a constant rate $r$ fixed every day for $n$ days ($d_i=1$), show that
      $R^{\mathrm{comp}}-r=\frac{(n-1)r^2}{720}+O(r^3)$. Evaluate for $r=5.2\%$, $n=91$ and compare with the example of
      \cref{ch06:sec-rates}.
\end{enumerate}
\end{exercise}

\begin{exercise}[not from the course]
(a) Show that a note paying compounded SOFR in arrears on $N$,
plus $N$ at maturity, is worth $N$ at every reset date when
discounting is on the SOFR curve.

(b) Deduce \eqref{ch06:eq-float} again.

(c) USD SOFR swaps pay two business days after the end of each
period. Write the value of one floating period with this payment
lag and explain why the telescoping is no longer exact.
\end{exercise}

\bigskip

\begin{exercise}[not from the course]
Using the quotes of \cref{ch06:ex-boot}:

(a) bootstrap $Z_1,Z_2,Z_3$ by hand;

(b) show that CDF gives $Z_4=\sqrt{Z_3Z_5}$, write the scalar
equation for $Z_5$ and solve it numerically;

(c) compute the one-year forward rates $f(k-1,k)$,
$k=1,\dots,5$, and the 4-year par rate;

(d) prove \cref{ch06:prop-cdf} directly from the daily compounding
definition, and show that the forward curve jumps at a knot unless
the two adjacent segments have equal slopes of $\ln Z$.
\end{exercise}

\bigskip

\begin{exercise}[not from the course]
\leavevmode
\begin{enumerate}[(a)]
  \item Explain, without formulas, why the convexity adjustment makes
      the futures rate \emph{higher} than the forward rate, and what
      would change if margin balances earned a fixed rate instead of the
      overnight rate.
  \item With $\sigma=0.8\%$ per year, evaluate
      $\frac12\sigma^2t_1t_2$ for 3-month contracts expiring in 2, 5 and
      10 years. From which maturity does it exceed a bid--offer of
      0.25\,bp?
\end{enumerate}
\end{exercise}
\par\medskip

\begin{exercise}[not from the course]
A company has a 5-year \$50M floating-rate loan at
SOFR$+150$\,bp and wants to fix its cost.
\begin{enumerate}[(a)]
  \item Should it pay or receive fixed on a swap?
  \item With the 5-year par rate at 3.40\%, what all-in fixed
      rate does it lock in (ignore day-count and payment-date
      differences)?
  \item One year later the 4-year par rate is 3.90\% and
      $A_4=3.60$. What is the swap's value to the company,
      and where does the offsetting loss appear?
\end{enumerate}
\end{exercise}

\begin{exercise}[not from the course]\label{ch06:ex-compadv}
Firm A can borrow fixed at 4.00\% or floating at SOFR$+0.30\%$;
firm B at 5.20\% fixed or SOFR$+1.00\%$. A wants floating, B fixed.
\begin{enumerate}[(a)]
  \item Show that A has an absolute advantage in both markets but a
      comparative advantage in fixed.
  \item Design a swap through a bank earning 0.10\% that splits the
      remaining gain equally between A and B.
  \item Which risk does each firm take on in exchange (think of
      \cref{ch06:sec-clearing})?
\end{enumerate}
\end{exercise}
\par\medskip

\begin{exercise}[not from the course]
\leavevmode
\begin{enumerate}[(a)]
  \item Show that the gamma of slide 31 equals
      $\epsilon\,\partial^2\mathrm{PV}/\partial r^2+O(\epsilon^3)$.
  \item For a 10-year zero-coupon bond with continuously compounded
      rate 4\% and notional \$100M, compute the PV01 and this gamma for
      $\epsilon=1$\,bp.
  \item How many revaluations does a full central-difference
      cross-gamma matrix need for $N$ curve inputs?
  \item Why does a Monte Carlo pricer need a larger $\epsilon$
      (\cref{ch:markets})?
\end{enumerate}
\end{exercise}

\begin{exercise}[not from the course]\label{ch06:ex-theta}
Annual-compounding zero rates are 3\% (1 year) and 4\% (2 years).
\begin{enumerate}[(a)]
  \item Price a 2-year zero-coupon bond.
  \item Value it one year later if the forwards are realised, and if
      the curve by tenor is unchanged.
  \item Split the one-year return in (b) into carry and roll-down,
      and relate each to the two theta conventions of
      \cref{ch06:sec-pnl}.
\end{enumerate}
\end{exercise}

\begin{exercise}[not from the course]
Using the bond of the example in \cref{ch06:sec-bonds} (DV01
$0.0454$ per 100) and the 5-year swap annuity $A_5=4.5178$:
\begin{enumerate}[(a)]
  \item which swap and which notional make \$10M face of the bond
      DV01-neutral?
  \item Show that for small $s$,
      $s\approx(P'_D-P_D)/\sum_it_i\mathrm{CF}_iZ(t_i)$, and check it
      against $s=25.6$\,bp.
  \item Which risk remains after the hedge?
\end{enumerate}
\end{exercise}

\begin{exercise}[not from the course]\label{ch06:ex-muni}
\leavevmode
\begin{enumerate}[(a)]
  \item At which tax rate is a 2.8\% muni equivalent to a 4\%
      taxable bond?
  \item A 10-year muni is priced to yield 70\% of a 4\%
      taxable yield. If investors come to expect the tax rate to
      fall from 30\% to 20\%, estimate the price change with a
      modified duration of 8, all else equal. Connect this to the
      Q\&A remark on tax policy.
\end{enumerate}
\end{exercise}
